\documentclass{haskellbeamer}

\title{Fixing Data-flow Problems in Syntax Trees}
\date{September 9, 2017}

\bibliography{bibliography.bib}
% Lattice

\newcommand{\llub}{\sqcup}
\newcommand{\lbiglub}{\bigsqcup}
\newcommand{\lleq}{\sqsubseteq}
\newcommand{\lless}{\sqsubset}
\newcommand{\lPair}[2]{\left\langle{}#1,#2\right\rangle{}}
\newcommand{\lTriple}[3]{\left\langle{}#1,#2,#3\right\rangle{}}
\newcommand{\both}{\,\&\,}
\newcommand{\mto}{\to_{+}}

% Domain

\newcommand{\sFVs}{\text{\textsf{FVs}}}
\newcommand{\sStr}{\text{\textsf{Str}}}
\newcommand{\sStrType}{\text{\textsf{StrType}}}
\newcommand{\sStrTrans}{\text{\textsf{StrTrans}}}

% Transfer function
\newcommand{\transfer}[2]{\mathcal{T}\llbracket#1\rrbracket_{#2}}
\newcommand*{\letupsc}{\textsc{LetUp}\xspace}
\newcommand*{\letdnsc}{\textsc{LetDn}\xspace}
\newcommand{\fix}{\keyword{fix}\xspace}
\newcommand{\up}{\sfop{up}}
\newcommand{\down}{\sfop{down}}

% Map Notation
\newcommand{\pfun}{\rightharpoonup}
\newcommand{\emptymap}{[]}
\newcommand{\id}[1]{#1}
\newcommand{\maplit}[3][\id]{\left[#1{#2\mapsto#3}\right]}
\newcommand{\restrict}[2]{#1\restriction_{#2}}


\begin{document}

\maketitle

\begin{frame}{Introduction}
  \begin{itemize}
    \item My master's thesis\footnote{\tiny\url{https://pp.ipd.kit.edu/uploads/publikationen/graf17masterarbeit.pdf}}: Call Arity vs. Demand Analysis
      \begin{itemize}
        \item Result: Usage Analysis generalising Call Arity
        \item Precision of Call Arity without co-call graphs
      \end{itemize}
    \item Requirements led to complex analysis order
    \item \emph{Specification} of data-flow problem decoupled from its \emph{solution}
  \end{itemize}
\end{frame}

\begin{frame}[fragile]{Strictness Analysis}
  \begin{itemize}
    \item Provides lower bounds on \emph{evaluation cardinality}
    \item Is this variable evaluated at least once?
      \begin{itemize}
        \item \emph{Strictness}: $\sStr \Coloneqq S \mid L$
        \item $S$trict (Yes!)
        \item $L$azy (Not sure)
      \end{itemize}
    \item Enables call-by-value, unboxing
  \end{itemize}
  \begin{center}
    \begin{minipage}{0.7\textwidth}
      \begin{overprint}
        \onslide<1>
        \begin{haskell}
          main = do
            let  x = ... -- $S$
            let  y = ... -- $S$
            let  z = ... -- $L$
            print (x + if odd y then y else z)
        \end{haskell}
        \onslide<2>
        \begin{haskell}
          main = do
            let !x = ... -- $S$
            let !y = ... -- $S$
            let  z = ... -- $L$
            print (x + if odd y then y else z)
        \end{haskell}
      \end{overprint}
    \end{minipage}
  \end{center}
\end{frame}

\begin{frame}[fragile]{GHC's Demand Analyser}
  \begin{itemize}
    \item Performs strictness analysis (among other things)
    \item Fuels Worker/Wrapper transformation
    \item Backward analysis
      \begin{itemize}
        \item Which strictness does an expression place on its free variables?
        \item Which strictness does a function place on its arguments?
      \end{itemize}
    \item \emph{Strictness type}: $\sStrType = \lPair{\sFVs \pfun \sStr}{\sStr^*}$
  \end{itemize}
\end{frame}

\begin{frame}[fragile]{Strictness Signatures}
  \begin{itemize}
    \item Looks at the right-hand side of \hs{const} before the \hs{let} body!
    \item \emph{Unleashes} strictness type of \hs{const}'s RHS at call sites
  \end{itemize}
  \begin{center}
    \begin{minipage}{0.8\textwidth}
      \begin{haskell*}{escapeinside=||}
        let const a b = a -- $\mathtt{const} :: \lPair{\emptymap}{[S, L]}$
        in const 
            y            -- $S$
            (error |\PYG{l+s}{\PYGZdq{}}\PYG{l+s}{\includegraphics[width=3.85mm]{bomb.png}}\PYG{l+s}{\PYGZdq{}}|) -- $L$
      \end{haskell*}
    \end{minipage}
  \end{center}
\end{frame}

\begin{frame}[fragile]{Call Context Matters}
  \begin{overprint}
    \onslide<1-3>
    \begin{itemize}
      \item Whole expression is strict in \hs{z}
      \item Only digests \hs{f} for manifest arity 1, can't look under lambda
      \item \hs{f} is called with 2 arguments
      \item<4>[\ding{213}] Analyse bound function when incoming arity is known
    \end{itemize}
    \onslide<4->
    \begin{itemize}
      \item Solution: Analyse RHS when incoming arity is known
      \item Formally: Finite approximation of \emph{strictness transformer}
        \begin{itemize}
          \item $\sStrTrans = \mathbb{N} \to \sStrType$
        \end{itemize}
      \item Exploit laziness to memoise results?
    \end{itemize}
  \end{overprint}
  \begin{overprint}
    \onslide<1>
    \begin{center}
      \begin{minipage}{0.7\textwidth}
        \begin{haskell*}{escapeinside=||}
          let f x = -- $\mathtt{f} :: \lPair{\maplit{\mathtt{z}}{L}}{[S]}$
                if odd x
                  then \y -> y*|\color{red}\texttt{z}|
                  else \y -> y+|\color{red}\texttt{z}|
          in f 1 2
        \end{haskell*}
      \end{minipage}
    \end{center}
    \onslide<2>
    \begin{center}
      \begin{minipage}{0.7\textwidth}
        \begin{haskell*}{escapeinside=||}
          let f x = -- $\mathtt{f} :: \lPair{\maplit{\mathtt{z}}{L}}{[S]}$
                if odd x
                  then \y -> y*z
                  else \y -> y+z
          in seq (f 1) 42
        \end{haskell*}
      \end{minipage}
    \end{center}
    \onslide<3>
    \begin{center}
      \begin{minipage}{0.7\textwidth}
        \begin{haskell*}{escapeinside=||}
          let f x = -- $\mathtt{f} :: \lPair{\maplit{\mathtt{z}}{L}}{[S]}$
                if odd x
                  then \y -> y*z
                  else \y -> y+z
          in f 1 2
        \end{haskell*}
      \end{minipage}
    \end{center}
    \onslide<4>
    \begin{center}
      \begin{minipage}{0.7\textwidth}
        \begin{haskell*}{escapeinside=||}
          let f x = -- $\mathtt{f}_{\color{red}1} :: \lPair{\maplit{\mathtt{z}}{L}}{[S]}$
                if odd x
                  then \y -> y*z
                  else \y -> y+z
          in f 1 2
        \end{haskell*}
      \end{minipage}
    \end{center}
    \onslide<5>
    \begin{center}
      \begin{minipage}{0.7\textwidth}
        \begin{haskell*}{escapeinside=||}
          let f x = -- $\mathtt{f}_{\color{red}2} :: \lPair{\maplit{\mathtt{z}}{S}}{[S,S]}$
                if odd x
                  then \y -> y*z
                  else \y -> y+z
          in f 1 2
        \end{haskell*}
      \end{minipage}
    \end{center}
  \end{overprint}
\end{frame}
  
\begin{frame}[fragile]{Recursion}
  \begin{itemize}
    \item Exploit laziness to memoise approximations?
    \item[\xmark] Recursion leads to termination problems
    \item Rediscovered fixed-point iteration, detached from the syntax tree
    \item Leads to data-flow problem, solved by worklist algorithm
  \end{itemize}
  \begin{center}
    \begin{minipage}{0.5\textwidth}
      \begin{haskell}
        let fac n = 
              if n == 0
                then 1
                else n * fac (n-1)
        in fac 12
      \end{haskell}
    \end{minipage}
  \end{center}
\end{frame}

\begin{frame}[fragile]{Data-flow Graph for Strictness Analysis}
  \begin{itemize}
    \item Allocate nodes to break recursion
      \begin{itemize}
        \item One top-level node
        \item One node per pair of $(\text{\hs{let} binding}, \text{incoming arity})$
      \end{itemize}
    \item Initialise worklist to top-level node
    \item Initialise nodes with $\bot$
  \end{itemize}
  \begin{overprint}
    \begin{columns}
      \begin{column}{0.35\textwidth}
        \begin{haskell}
          let f 0 = const 0
              f 1 = id 
              f n = f (n `mod` 2)
          in f x y
        \end{haskell}
      \end{column}
      \begin{column}{0.02\textwidth}
        {\Huge$\Rightarrow$}
      \end{column}
      \begin{column}{0.4\textwidth}
        \begin{tikzpicture}[->, thick]
          \useasboundingbox (-1,-1.5) rectangle (4,1.5);
          \node[n] at (0,0) (root) 
            [label={[visible on=<4>]above:$\bot$}]
            {$\dfnode{<root>}{0}$};
          \foreach \i in {-0.05,0,0.05} {
            \node[n, visible on=<2->] at (3,0) (fx) 
              [xshift=\i cm, yshift=-\i cm] {$\dfnode{f}{n}$};
          }
          \node[visible on=<4>] at ([shift={(90:1.1)}]fx) {$\bot$};
        \end{tikzpicture}
      \end{column}
    \end{columns}
    \onslide<3->
    \begin{columns}
      \begin{column}{0.55\textwidth}
        \hfill Worklist:
      \end{column}
      \begin{column}{0.55\textwidth}
        $\{\dfnode{<root>}{0}\}$ \hfill
      \end{column}
    \end{columns}
  \end{overprint}
\end{frame}

\begin{frame}[fragile]{Data-flow Graph for Strictness Analysis}
  \begin{overprint}
    \begin{columns}
      \begin{column}{0.35\textwidth}
        \begin{haskell}
          let f 0 = const 0
              f 1 = id 
              f n = f (n `mod` 2)
          in f x y
        \end{haskell}
      \end{column}
      \begin{column}{0.02\textwidth}
        {\Huge$\Rightarrow$}
      \end{column}
      \begin{column}{0.4\textwidth}
        \begin{tikzpicture}[->, thick]
          \useasboundingbox (-1,-2) rectangle (4,2);
          \node[n] at (0,0) (root) [hl={<2-8>}] 
            [label={[visible on=<-7>]above:$\bot$},
             label={[visible on=<8->]above:$\lPair{\alert<8,13>{\left[\mathtt{x} \mapsto S, \mathtt{y} \mapsto L\right]}}{[]}$}]
            {$\dfnode{<root>}{0}$};
          \foreach \i in {-0.05,0,0.05} {
            \node<-2>[n] at (3,0) (fx) [xshift=\i cm, yshift=-\i cm] {$\dfnode{f}{n}$};
          }
          \draw<3-> (root) edge [hl={<3,7>}] (f2);
          \node[n, visible on=<3->, hl={<3-7,10-12>}] 
            at (3,0) (f2) 
            [label={[visible on=<-4>]above:$\bot$},
             label={[visible on=<5-6>]above:$\bot=\lPair{\emptymap}{[S,S]}$},
             label={[visible on=<7->]above:$\lPair{\emptymap}{[S,\alert<7>{L}]}$}]
            {$\dfnode{f}{2}$};
          \draw<4-> (f2) edge [loop, in=-110, out=-70, looseness=5, hl={<4-6,11>}] (f2);
        \end{tikzpicture}
      \end{column}
    \end{columns}
    \begin{columns}
      \begin{column}{0.55\textwidth}
        \hfill Worklist:
      \end{column}
      \begin{column}{0.55\textwidth}
        \only<1>{$\{\alert{\dfnode{<root>}{0}}\}$}\only<2-5>{$\{\}$}\only<6-9>{$\{\alert<6,9>{\dfnode{f}{2}}\}$}\only<10->{$\{\}$}
      \end{column}
    \end{columns}
  \end{overprint}
\end{frame}

\begin{frame}[fragile]{Implementation}
  \begin{itemize}
    \item Hide iteration strategy behind \hs{TransferFunction} monad
    \item Data-flow nodes \hs{k}, denoting lattice \hs{v}
    \item Single `impure' primitive \hs{dependOn}
  \end{itemize}
  \begin{center}
    \begin{minipage}{0.9\textwidth}
      \begin{haskell}
        data TransferFunction k v a
        instance Monad (TransferFunction k v)

        dependOn
          :: Ord k
          => k
          -> TransferFunction k v (Maybe v)
      \end{haskell}
    \end{minipage}
  \end{center}
\end{frame}

\begin{frame}[fragile]{Implementation}
  \begin{itemize}
    \item \hs{DataFlowProblem} assigns \hs{TransferFunction} and \hs{ChangeDetector} to nodes
  \end{itemize}
  \begin{center}
    \begin{minipage}{0.9\textwidth}
      \begin{haskell}
        type ChangeDetector k v
          = k -> v -> v -> Bool

        data DataFlowProblem k v 
          = DFP 
          { transfer      :: k -> TransferFunction k v v
          , detectChanges :: k -> ChangeDetector k v
          }
      \end{haskell}
    \end{minipage}
  \end{center}
\end{frame}

\begin{frame}[fragile]{Implementation}
  \begin{itemize}
    \item \hs{fixProblem} solves data-flow problems 
    \item Specification as \hs{DataFlowProblem}
    \item Implements fixed-point iteration strategy
      \begin{itemize}
        \item Can use worklist algorithm, starting from a specified root set
      \end{itemize}
  \end{itemize}
  \begin{center}
    \begin{minipage}{0.5\textwidth}
      \begin{haskell}
        fixProblem
          :: Ord k
          => DataFlowProblem k v
          -> Set k
          -> Map k v
      \end{haskell}
    \end{minipage}
  \end{center}
\end{frame}

\begin{frame}[fragile]{Applied to Strictness Analysis}
  \begin{itemize}
    \item Denote expressions by their strictness transformer
    \item Model points of strictness transformer separately
    \item Instantiate as \hs{DataFlowProblem (ExprNode, Arity) StrType}
    \item \hs{ExprNode}: Totally ordered, allocated as needed
      \begin{itemize}
        \item Dictates priority in worklist
        \item Performance depends on suitable priorities
      \end{itemize}
  \end{itemize}
\end{frame}

\begin{frame}{Comparison to \texttt{hoopl}}
  \begin{itemize}
    \item \texttt{hoopl} \parencite{hoopl} works on CFGs
      \begin{itemize}
        \item Our data-flow graph is much less restrictive
        \item Edges implicit in DSL 
      \end{itemize}
    \item Designed for imperative languages
    \item `Operational' rather than `denotational'
      \begin{itemize}
        \item Analysis along control-flow rather than data-flow
      \end{itemize}
    \item Makes (join-semi)lattice explicit
      \begin{itemize}
        \item TODO
      \end{itemize}
    \item Also includes a solution for transformations
  \end{itemize}
\end{frame}

\begin{frame}{Discussion}
  \begin{itemize}
    \item[\cmark] Decouple analysis logic from iteration logic by a graph-based approach
    \item[\xmark] Coupling not as painful as it would be in imperative programs
    \item[\cmark] Still obscures intent, even obstructs ideas
    \item[\cmark] `Hacks' such as caching of analysis results as in \textcite[§9.2]{dmd} between iterations for free
    \item[\xmark] Unclear how performance is affected
    \item[\xmark] Can only shine if shared concerns are actually extracted from a number of analyses
  \end{itemize}
\end{frame}

\begin{frame}{Conclusion}
  \begin{itemize}
    \item Pitched an interesting idea that came out of my thesis
    \item Separate \emph{specification} of the data-flow problem from computing its \emph{solution}
    \item Graph-based abstraction, single solver
    \item Future Work:
      \begin{enumerate}
        \item Monotone maps with partially-ordered keys\footnote{\url{https://github.com/sgraf812/pomaps/}}
        \item Polish API, make a package
        \item Testdrive and measure it in GHC
      \end{enumerate}
  \end{itemize}
\end{frame}

\begin{frame}{Bibliography} 
  \printbibliography[heading=none]
\end{frame}

\section{Backup}

\begin{frame}[fragile]{Example}
  \begin{overprint}
    \begin{columns}
      \begin{column}{0.4\textwidth}
        \begin{overprint}
          \onslide<1>
          \begin{haskell}
            let f 0 = const 0
                f 1 = id 
                f n = 
                  const (f (n `mod` 2) 4)
            in f x y
          \end{haskell}
          \onslide<2->
          \begin{haskell}
            let f 0 = const 0
                f 1 = id 
                f n = 
                  const (f (n `mod` 2) 4)
            in seq (f x) (f x y)
          \end{haskell}
        \end{overprint}
      \end{column}
      \begin{column}{0.02\textwidth}
        {\Huge$\Rightarrow$}
      \end{column}
      \begin{column}{0.4\textwidth}
        \begin{tikzpicture}[->, thick]
          \useasboundingbox (-1,-2) rectangle (4,2);
          \node[n] at (0,0) (root) {$\dfnode{<root>}{0}$};
          \foreach \i in {-0.05,0,0.05} {
            \node<-2>[n] at (3,0) (fx) [xshift=\i cm, yshift=-\i cm] {$\dfnode{f}{n}$};
          }
          \node[n, visible on=<3->] at (3,1.5) (f1) {$\dfnode{f}{1}$};
          \node[n, visible on=<3->] at (3,-1.5) (f2) {$\dfnode{f}{2}$};
          \draw<3-> 
            (root) edge [bend left] (f1)
            (root) edge [bend right] (f2)
            (f1) edge [bend left] (f2)
            (f2) edge [loop, in=-80, out=-40, looseness=5] (f2);
        \end{tikzpicture}
      \end{column}
    \end{columns}
  \end{overprint}
\end{frame}

\begin{frame}[fragile]{Implementation: Behind the Curtain}
  \begin{itemize}
    \item \hs{TransferFunction} is a \hs{State} monad around \hs{WorklistState}
  \end{itemize}
  \begin{center}
    \begin{minipage}{0.8\textwidth}
      \begin{haskell}
        data TransferFunction node lattice a
          = TFM (State (WorklistState node lattice ) a)
          deriving (Functor, Applicative, Monad)
      \end{haskell}
    \end{minipage}
  \end{center}
\end{frame}

\begin{frame}[fragile]{Threading annotated expressions}
  \begin{itemize}
    \item Annotated \hs{CoreExpr}s are the reason why we do this!
    \item Thread it through all nodes: \hs{DataFlowProblem (ExprNode, Arity) (StrType, CoreExpr)}
    \item Complicates change detection
      \begin{itemize}
        \item Expressions follow AST structure
        \item Possibly change when strictness type did not
        \item \hs{ChangeDetector} has to check set of changed dependencies
      \end{itemize}
    \item $\sStr \Coloneqq S \mid L$ not enough for annotating functions
      \begin{itemize}
        \item $\sStr \Coloneqq S^n \mid L$ with arity $n\in \mathbb{N}$
        \item `\hs{f} was called at least once, with at least $n$ arguments'
      \end{itemize}
    \item ... Or do it as the Demand Analyser does: Assume manifest arity for annotation
      \begin{itemize}
        \item Be careful not to inline unsaturated wrappers!
      \end{itemize}
  \end{itemize}
\end{frame}

\begin{frame}[fragile]{Caching of Analysis Results due to Henglein}
  \begin{center}
    \begin{minipage}{0.8\textwidth}
      \begin{haskell}
        let f x =
              let g y = 
                    if odd y 
                      then g (y - 1) 
                      else x
              in if even x 
                  then g x 
                  else f (3*x + 1)
        in f 7
      \end{haskell}
    \end{minipage}
  \end{center}
\end{frame}

\end{document}